%!TeX root=../tese.tex
%("dica" para o editor de texto: este arquivo é parte de um documento maior)
% para saber mais: https://tex.stackexchange.com/q/78101

% Vamos definir alguns comandos auxiliares para facilitar.

% "textbackslash" é muito comprido.
\newcommand{\sla}{\textbackslash}

% Vamos escrever comandos (como "make" ou "itemize") com formatação especial.
\newcommand{\cmd}[1]{\textsf{#1}}

% Idem para packages; aqui estamos usando a mesma formatação de \cmd,
% mas poderíamos escolher outra.
\newcommand{\pkg}[1]{\textsf{#1}}

% A maioria dos comandos LaTeX começa com "\"; vamos criar um
% comando que já coloca essa barra e formata com "\cmd".
\newcommand{\ltxcmd}[1]{\cmd{\sla{}#1}}

\chapter{Resultados e Discussão}
\label{chap:resultados}

\section{Concordância Inter avaliadores na classificação da relação entre e-mails e commits}
A classificação independente abrangeu 1.292 pares formados por
\textit{e-mails} de submissão de patches e seus possíveis \textit{commits}
correspondentes. Foram excluídos os \textit{commits} de merge, pois eles
reúnem múltiplos \textit{commits} em uma única integração e poderiam
comprometer a unidade de análise. Nesse conjunto, os avaliadores
concordaram em 1.234 casos, equivalentes a 95,51\% das avaliações
válidas, e divergiram em 58 casos (4,49\%). A Tabela 1 sintetiza esses
indicadores globais de concordância.

\textbf{Tabela 1 -- Indicadores globais de concordância Inter
	avaliadores}	
\begin{longtable}[]{@{}
			>{\raggedright\arraybackslash}p{(\linewidth - 2\tabcolsep) * \real{0.5659}}
			>{\centering\arraybackslash}p{(\linewidth - 2\tabcolsep) * \real{0.3144}}@{}}
		\toprule\noalign{}
		\begin{minipage}[b]{\linewidth}\raggedright
			\textbf{Indicador}
		\end{minipage} & \begin{minipage}[b]{\linewidth}\centering
			\textbf{Resultado}
		\end{minipage} \\
		\midrule\noalign{}
		\endhead
		\bottomrule\noalign{}
		\endlastfoot
		Avaliações válidas & 1.292 \\
		Concordâncias & 1.234 \\
		Desacordos & 58 \\
		Concordância percentual & 95,51\% \\
		Kappa de Cohen & 0,8219 \\
\end{longtable}

Embora a concordância percentual observada tenha sido elevada, essa
medida não desconta o acordo esperado ao acaso. Por esse motivo, a
análise foi complementada pelo \textit{kappa} de Cohen. O valor $\kappa$ = 0,8219
situa-se, segundo a classificação convencional de \cite{landis1977measurement},
no intervalo de 0,81 a 1,00, usualmente denominado concordância ``quase
perfeita''. Assim, a combinação entre 95,51\% de acordo observado e
\textit{kappa} superior a 0,80 fornece evidência de elevada
reprodutibilidade do procedimento manual.
%\newpage
\textbf{Tabela 2 -- Matriz de concordância entre os avaliadores}
\begin{longtable}[]{@{}
			>{\raggedright\arraybackslash}p{(\linewidth - 6\tabcolsep) * \real{0.3668}}
			>{\centering\arraybackslash}p{(\linewidth - 6\tabcolsep) * \real{0.2315}}
			>{\centering\arraybackslash}p{(\linewidth - 6\tabcolsep) * \real{0.2315}}
			>{\centering\arraybackslash}p{(\linewidth - 6\tabcolsep) * \real{0.1703}}@{}}
		\toprule\noalign{}
		\begin{minipage}[b]{\linewidth}\centering
			\textbf{Classificação do Avaliador 1}
		\end{minipage} & \begin{minipage}[b]{\linewidth}\centering
			\textbf{Avaliador 2: Não}
		\end{minipage} & \begin{minipage}[b]{\linewidth}\centering
			\textbf{Avaliador 2: Sim}
		\end{minipage} & \begin{minipage}[b]{\linewidth}\centering
			\textbf{Total}
		\end{minipage} \\
		\midrule\noalign{}
		\endhead
		\bottomrule\noalign{}
		\endlastfoot
		Avaliador 1: Não & 162 & 21 & 183 \\
		Avaliador 1: Sim & 37 & 1.072 & 1.109 \\
		Total & 199 & 1.093 & 1.292 \\
\end{longtable}

A Tabela 2 detalha a composição desse acordo. A diagonal principal
concentra os 1.234 julgamentos concordantes: 162 casos classificados
como ``Não'' por ambos os avaliadores e 1.072 classificados como
``Sim''. Entre os 58 desacordos, o Avaliador 1 registrou ``Sim'' e o
Avaliador 2, ``Não'' em 37 casos; nos 21 restantes, ocorreu o padrão
inverso. Consequentemente, o Avaliador 1 atribuiu ``Sim'' em 85,84\% dos
casos, ante 84,60\% do Avaliador 2. A diferença marginal, de apenas 1,24
ponto percentual, sugere leve maior propensão do Avaliador 1 à categoria
``Sim'', mas sua pequena magnitude deve ser interpretada em conjunto com
a elevada concordância global.

\textbf{Tabela 3 -- Distribuição das categorias por avaliador}
\begin{longtable}[]{@{}
		>{\raggedright\arraybackslash}p{(\linewidth - 6\tabcolsep) * \real{0.3668}}
		>{\centering\arraybackslash}p{(\linewidth - 6\tabcolsep) * \real{0.2315}}
		>{\centering\arraybackslash}p{(\linewidth - 6\tabcolsep) * \real{0.2315}}
		>{\centering\arraybackslash}p{(\linewidth - 6\tabcolsep) * \real{0.1703}}@{}}
	% conteúdo da tabela
	\toprule\noalign{}
	\begin{minipage}[b]{\linewidth}\centering
		\textbf{Classificação do Avaliador 1}
	\end{minipage} & \begin{minipage}[b]{\linewidth}\centering
		\textbf{Avaliador 2: Não}
	\end{minipage} & \begin{minipage}[b]{\linewidth}\centering
		\textbf{Avaliador 2: Sim}
	\end{minipage} & \begin{minipage}[b]{\linewidth}\centering
		\textbf{Total}
	\end{minipage} \\
	\midrule\noalign{}
	\endhead
	\bottomrule\noalign{}
	\endlastfoot
	Avaliador 1: Não & 162 & 21 & 183 \\
	Avaliador 1: Sim & 37 & 1.072 & 1.109 \\
	Total & 199 & 1.093 & 1.292 \\
	
\end{longtable}

Como evidencia a Tabela 3, a predominância da categoria ``Sim'' foi
consistente entre os avaliadores. Esse desequilíbrio merece atenção
porque a concordância percentual pode ser inflada quando uma categoria é
muito frequente. O \textit{kappa} de Cohen atenua parcialmente esse efeito
ao considerar o acordo esperado a partir das distribuições marginais.
Ainda assim, a prevalência precisa permanecer explícita na
interpretação. Por essa razão, a Tabela 4 apresenta separadamente os
acordos específicos por categoria.

\textbf{Tabela 4 -- Concordância específica por categoria}
\begin{longtable}[]{@{}
			>{\raggedright\arraybackslash}p{(\linewidth - 8\tabcolsep) * \real{0.1645}}
			>{\centering\arraybackslash}p{(\linewidth - 8\tabcolsep) * \real{0.2080}}
			>{\centering\arraybackslash}p{(\linewidth - 8\tabcolsep) * \real{0.2238}}
			>{\centering\arraybackslash}p{(\linewidth - 8\tabcolsep) * \real{0.2238}}
			>{\centering\arraybackslash}p{(\linewidth - 8\tabcolsep) * \real{0.1799}}@{}}
		\toprule\noalign{}
		\begin{minipage}[b]{\linewidth}\centering
			\textbf{Categoria}
		\end{minipage} & \begin{minipage}[b]{\linewidth}\centering
			\textbf{Concordâncias}
		\end{minipage} & \begin{minipage}[b]{\linewidth}\centering
			\textbf{Atribuições do Avaliador 1}
		\end{minipage} & \begin{minipage}[b]{\linewidth}\centering
			\textbf{Atribuições do Avaliador 2}
		\end{minipage} & \begin{minipage}[b]{\linewidth}\centering
			\textbf{Acordo específico}
		\end{minipage} \\
		\midrule\noalign{}
		\endhead
		\bottomrule\noalign{}
		\endlastfoot
		Não & 162 & 183 & 199 & 84,82\% \\
		Sim & 1.072 & 1.109 & 1.093 & 97,37\% \\
\end{longtable}

A concordância específica foi maior para ``Sim'' (97,37\%) do que para
``Não'' (84,82\%). Portanto, os avaliadores reconheceram com mais
uniformidade os casos em que o e-mail estava associado ao \textit{commit},
enquanto a categoria menos frequente concentrou proporcionalmente mais
incerteza. Embora o acordo de 84,82\% para ``Não'' ainda seja elevado, a
diferença de 12,55 pontos percentuais entre as categorias recomenda
examinar os 58 desacordos, sobretudo em relação a critérios
potencialmente ambíguos, como alterações no assunto, ausência de índices
do \textit{diff}, divergências de autoria e intervalos temporais
limítrofes.

Em conjunto, esses resultados indicam que o protocolo de classificação
independente foi aplicado de modo consistente. A elevada concordância
global, o \textit{kappa} de 0,8219 e os acordos específicos superiores a
84\% sustentam o uso da classificação manual como referência
suficientemente estável para a etapa de modelagem. Entretanto,
concordância não equivale a validade, pois dois avaliadores podem
aplicar de forma semelhante um critério inadequado. Para reduzir esse
risco, a resolução consensual dos desacordos incluiu nova leitura do
texto integral do e-mail e das informações do \textit{commit}, com atenção
especial às datas, aos autores, ao código e aos arquivos alterados.

\section{Modelo logístico inicial e seleção das variáveis preditoras}
O modelo logístico inicial foi ajustado com 1.292 observações e oito
preditores. O ajuste convergiu, apresentou teste da razão de
verossimilhanças globalmente significativo (p \textless{} 0,001) e
pseudo-R² de McFadden igual a 0,8228. Esses resultados indicam que o
conjunto inicial de preditores acrescentou informação substancial em
comparação com o modelo que continha apenas o intercepto. Contudo, um
bom ajuste global não assegura contribuição independente de cada
coeficiente, sobretudo quando há forte associação entre as variáveis
explicativas. A Tabela 5 resume o ajuste inicial.

\textbf{Tabela 5 -- Síntese do ajuste logístico inicial}
\begin{longtable}[]{@{}
			>{\raggedright\arraybackslash}p{(\linewidth - 2\tabcolsep) * \real{0.5039}}
			>{\raggedright\arraybackslash}p{(\linewidth - 2\tabcolsep) * \real{0.2975}}@{}}
		\toprule\noalign{}
		\begin{minipage}[b]{\linewidth}\raggedright
			\textbf{Indicador}
		\end{minipage} & \begin{minipage}[b]{\linewidth}\raggedright
			\textbf{Resultado}
		\end{minipage} \\
		\midrule\noalign{}
		\endhead
		\bottomrule\noalign{}
		\endlastfoot
		Observações & 1.292 \\
		Registros removidos por ausência & 0 \\
		Preditores & 8 \\
		Log-verossimilhança & -92,452 \\
		Pseudo-R² de McFadden & 0,8228 \\
		p do teste da razão de verossimilhanças & \textless{} 0,001 \\
		Convergência & Sim \\
\end{longtable}


Para investigar possíveis sobreposições entre os preditores, examinou-se
primeiro a matriz de correlação. Adotou-se, como critério descritivo,
correlação absoluta igual ou superior a 0,80. Nove pares atenderam a
esse critério, com valores entre 0,8225 e 0,9716. A associação mais
forte ocorreu entre CoincideAutor e Intervalo (r = 0,9716), seguida das
relações entre CoincideAutor e CossenoDiffs (r = 0,9417), Intervalo e
CossenoDiffs (r = 0,9303) e CossenoDiffs e SimilaridadeDiffs (r =
0,9284), conforme a Tabela 6.

\newpage

\textbf{Tabela 6 -- Pares de preditores com correlação absoluta igual ou
	superior a 0,80}
\begin{longtable}[]{@{}
			>{\raggedright\arraybackslash}p{(\linewidth - 4\tabcolsep) * \real{0.3357}}
			>{\raggedright\arraybackslash}p{(\linewidth - 4\tabcolsep) * \real{0.3368}}
			>{\raggedright\arraybackslash}p{(\linewidth - 4\tabcolsep) * \real{0.3274}}@{}}
		\toprule\noalign{}
		\begin{minipage}[b]{\linewidth}\raggedright
			\textbf{Variável 1}
		\end{minipage} & \begin{minipage}[b]{\linewidth}\raggedright
			\textbf{Variável 2}
		\end{minipage} & \begin{minipage}[b]{\linewidth}\raggedright
			\textbf{Correlação}
		\end{minipage} \\
		\midrule\noalign{}
		\endhead
		\bottomrule\noalign{}
		\endlastfoot
		CoincideAutor & Intervalo & 0,9716 \\
		CoincideAutor & CossenoDiffs & 0,9417 \\
		Intervalo & CossenoDiffs & 0,9303 \\
		CossenoDiffs & SimilaridadeDiffs & 0,9284 \\
		CoincideAssunto & CoincideAutor & 0,8679 \\
		CoincideAssunto & CossenoDiffs & 0,8649 \\
		CoincideAssunto & Intervalo & 0,8616 \\
		CoincideAutor & SimilaridadeDiffs & 0,8412 \\
		Intervalo & SimilaridadeDiffs & 0,8225 \\
\end{longtable}

Como as correlações bivariadas não resumem toda a dependência conjunta
entre os preditores, o diagnóstico foi complementado pelo fator de
inflação da variância (VIF). Essa medida quantifica quanto a variância
estimada de um coeficiente é ampliada pela associação da variável com os
demais preditores. O VIF foi empregado como diagnóstico contextual, e
não como regra automática de exclusão. \cite{obrien2007caution}  O'Brien (2007) adverte que
limiares usuais, como VIF igual a 10, não devem determinar isoladamente
a eliminação de variáveis; a decisão também deve considerar o objetivo
do modelo, a redundância substantiva, a precisão dos coeficientes e a
estabilidade do ajuste.

\textbf{Tabela 7 -- Fator de inflação da variância no modelo inicial}
\begin{longtable}[]{@{}
			>{\raggedright\arraybackslash}p{(\linewidth - 6\tabcolsep) * \real{0.2939}}
			>{\raggedright\arraybackslash}p{(\linewidth - 6\tabcolsep) * \real{0.1467}}
			>{\raggedright\arraybackslash}p{(\linewidth - 6\tabcolsep) * \real{0.1891}}
			>{\raggedright\arraybackslash}p{(\linewidth - 6\tabcolsep) * \real{0.2679}}@{}}
		\toprule\noalign{}
		\begin{minipage}[b]{\linewidth}\raggedright
			\textbf{Variável}
		\end{minipage} & \begin{minipage}[b]{\linewidth}\raggedright
			\textbf{VIF}
		\end{minipage} & \begin{minipage}[b]{\linewidth}\raggedright
			\textbf{Tolerância}
		\end{minipage} & \begin{minipage}[b]{\linewidth}\raggedright
			\textbf{Diagnóstico inicial}
		\end{minipage} \\
		\midrule\noalign{}
		\endhead
		\bottomrule\noalign{}
		\endlastfoot
		CossenoDiffs & 23,0786 & 0,0433 & Elevada \\
		CoincideAutor & 22,7139 & 0,0440 & Elevada \\
		Intervalo & 19,7402 & 0,0507 & Elevada \\
		SimilaridadeDiffs & 8,1359 & 0,1229 & Moderada \\
		CoincideAssunto & 4,6628 & 0,2145 & Limítrofe \\
		CoincideIndices & 1,2013 & 0,8324 & Baixa \\
		CoincideLink & 1,1632 & 0,8597 & Baixa \\
		Delta & 1,0863 & 0,9206 & Baixa \\
\end{longtable}


Os resultados da Tabela 7 confirmam multicolinearidade importante no
bloco formado por CossenoDiffs, CoincideAutor, Intervalo,
SimilaridadeDiffs e CoincideAssunto. Os três primeiros apresentaram VIFs
de 23,0786, 22,7139 e 19,7402, respectivamente, acompanhados de
tolerâncias próximas de 0,05. SimilaridadeDiffs apresentou VIF de
8,1359, enquanto CoincideAssunto alcançou 4,6628. Conforme discutem \cite{dormann2013collinearity}, correlações elevadas entre preditores podem
inflar a variância dos parâmetros, dificultar a identificação dos
efeitos relevantes e comprometer a transferência do modelo quando a
estrutura de correlação se altera. Diante desse diagnóstico, os
coeficientes do modelo inicial, apresentados na Tabela 8, foram
examinados em conjunto com os indicadores de colinearidade.

\newpage

\textbf{Tabela 8 -- Coeficientes do modelo inicial relacionados à
	decisão de seleção}

\begin{longtable}[]{@{}
			>{\raggedright\arraybackslash}p{(\linewidth - 8\tabcolsep) * \real{0.2150}}
			>{\raggedright\arraybackslash}p{(\linewidth - 8\tabcolsep) * \real{0.2040}}
			>{\raggedright\arraybackslash}p{(\linewidth - 8\tabcolsep) * \real{0.1961}}
			>{\raggedright\arraybackslash}p{(\linewidth - 8\tabcolsep) * \real{0.1925}}
			>{\raggedright\arraybackslash}p{(\linewidth - 8\tabcolsep) * \real{0.1924}}@{}}
		\toprule\noalign{}
		\begin{minipage}[b]{\linewidth}\raggedright
			\textbf{Variável}
		\end{minipage} & \begin{minipage}[b]{\linewidth}\raggedright
			\textbf{Coeficiente}
		\end{minipage} & \begin{minipage}[b]{\linewidth}\raggedright
			\textbf{Erro-padrão}
		\end{minipage} & \begin{minipage}[b]{\linewidth}\raggedright
			\textbf{p-value}
		\end{minipage} & \begin{minipage}[b]{\linewidth}\raggedright
			\textbf{IC 95\%}
		\end{minipage} \\
		\midrule\noalign{}
		\endhead
		\bottomrule\noalign{}
		\endlastfoot
		CoincideLink & 3,5658 & 1,241 & 0,004 & 1,133 a 5,999 \\
		CoincideAssunto & -0,0549 & 1,337 & \textit{\textbf{0,967}} & -2,676 a
		2,566 \\
		CoincideAutor & 4,6911 & 1,597 & 0,003 & 1,560 a 7,822 \\
		Intervalo & 2,1723 & 1,377 & \textit{\textbf{0,115}} & -0,526 a 4,870 \\
		CoincideIndices & 1,2247 & 0,528 & 0,020 & 0,190 a 2,260 \\
		Delta & -0,0488 & 0,157 & 0,756 & -0,356 a 0,259 \\
		CossenoDiffs & 0,2945 & 0,566 & \textit{\textbf{0,603}} & -0,815 a
		1,404 \\
		SimilaridadeDiffs & 1,2154 & 0,843 & 0,149 & -0,436 a 2,867 \\
\end{longtable}


A seleção das variáveis decorreu, portanto, da análise conjunta das
Tabelas 7 e 8. CoincideAssunto foi retirada porque sua informação estava
fortemente sobreposta à fornecida por CoincideAutor, CossenoDiffs e
Intervalo, com correlações de 0,8679, 0,8649 e 0,8616. Embora o VIF de
4,6628 não caracterizasse, isoladamente, multicolinearidade grave, seu
coeficiente foi praticamente nulo (-0,0549), o p-valor foi 0,967 e o
intervalo de confiança foi amplo (-2,676 a 2,566). Em conjunto, esses
elementos indicaram ausência de contribuição incremental identificável
no modelo completo.

CossenoDiffs, por sua vez, foi removida porque apresentou o maior VIF do
modelo (23,0786), além de forte correlação com SimilaridadeDiffs,
CoincideAutor, Intervalo e CoincideAssunto. Seu coeficiente também não
diferiu estatisticamente de zero (p = 0,603). Como CossenoDiffs e
SimilaridadeDiffs representam medidas de proximidade entre os conteúdos
comparados, a manutenção simultânea aumentava a redundância, ampliava o
erro-padrão e dificultava a atribuição de efeito independente. Optou-se,
assim, por preservar SimilaridadeDiffs como medida contínua de
similaridade textual e excluir CossenoDiffs.

Intervalo também foi excluída. A variável apresentou VIF de 19,7402 e
correlações elevadas com CoincideAutor, CossenoDiffs, CoincideAssunto e
SimilaridadeDiffs. Apesar do coeficiente positivo, o intervalo de
confiança incluía zero e o p-valor foi 0,115. Sua retirada reduziu a
sobreposição entre indicadores sem eliminar do modelo as dimensões de
autoria e similaridade que se pretendia representar.

É importante ressaltar que a seleção não se baseou apenas nos p-valores
nem em um limiar rígido de VIF. Foram considerados conjuntamente as
correlações cruzadas, a tolerância, a redundância conceitual, a
amplitude dos intervalos de confiança e o objetivo de obter um modelo
parcimonioso e interpretável. Essa estratégia evita decisões mecânicas
baseadas exclusivamente em regras de corte e reduz o risco de atribuir
significado substantivo a coeficientes instáveis produzidos pelo
compartilhamento de informação entre preditores correlacionados. A
Tabela 9 consolida as decisões adotadas.

\textbf{Tabela 9 -- Decisão sobre as variáveis do modelo inicial}
\begin{longtable}[]{@{}
			>{\raggedright\arraybackslash}p{(\linewidth - 4\tabcolsep) * \real{0.2349}}
			>{\raggedright\arraybackslash}p{(\linewidth - 4\tabcolsep) * \real{0.1970}}
			>{\raggedright\arraybackslash}p{(\linewidth - 4\tabcolsep) * \real{0.5681}}@{}}
		\toprule\noalign{}
		\begin{minipage}[b]{\linewidth}\raggedright
			\textbf{Variável}
		\end{minipage} & \begin{minipage}[b]{\linewidth}\raggedright
			\textbf{Decisão}
		\end{minipage} & \begin{minipage}[b]{\linewidth}\raggedright
			\textbf{Fundamento principal}
		\end{minipage} \\
		\midrule\noalign{}
		\endhead
		\bottomrule\noalign{}
		\endlastfoot
		CoincideLink & Mantida & Baixo VIF; associação independente no modelo
		inicial \\
		CoincideAssunto & Retirada & Redundância cruzada; coeficiente nulo; p =
		0,967 \\
		CoincideAutor & Mantida provisoriamente & Relevância substantiva e p =
		0,003; VIF deve ser recalculado \\
		Intervalo & Retirada & VIF elevado e sobreposição com autoria e
		similaridade \\
		CoincideIndices & Mantida & Baixo VIF e p = 0,020 \\
		Delta & Mantida para avaliação final & Baixo VIF; contribuição preditiva
		deve ser verificada fora da amostra \\
		CossenoDiffs & Retirada & Maior VIF e redundância com
		SimilaridadeDiffs \\
		SimilaridadeDiffs & Mantida & Preserva medida contínua de similaridade
		textual \\
\end{longtable}

Após as exclusões, o conjunto reduzido foi formado por CoincideLink,
CoincideAutor, CoincideIndices, Delta e SimilaridadeDiffs. A manutenção
de CoincideAutor deve ser descrita como provisória até que o VIF seja
recalculado no modelo reduzido, pois seu valor inicial de 22,7139 também
foi elevado. O VIF depende das demais variáveis incluídas e pode
diminuir substancialmente após a retirada de Intervalo, CossenoDiffs e
CoincideAssunto. A versão final do artigo deve apresentar os VIFs
recalculados e comparar o modelo inicial e o reduzido quanto à
estabilidade dos coeficientes, discriminação e calibração em dados não
utilizados no ajuste.

Em síntese, o modelo inicial apresentou ajuste global elevado, mas
continha preditores que descreviam aspectos parcialmente coincidentes da
correspondência entre e-mails e \textit{commits}. A retirada de
CoincideAssunto, CossenoDiffs e Intervalo buscou reduzir a inflação da
variância, evitar duplicação de informação e favorecer uma interpretação
mais estável dos efeitos mantidos. A escolha do modelo final deve ser
confirmada pelo diagnóstico de multicolinearidade após a redução e pelo
desempenho fora da amostra, especialmente pelas métricas de
discriminação e calibração probabilística.

\section{Particionamento da amostra, diagnóstico do modelo reduzido e ajuste no treinamento}
A amostra analítica de 1.292 pares de e-mail e \textit{commit} foi dividida de forma agrupada, de modo que observações pertencentes ao mesmo agrupamento lógico (\textit{commit}) não fossem distribuídas entre treinamento, validação e teste. O conjunto de treinamento reuniu 902 pares (69,81\%), a validação 194 pares (15,02\%) e o teste 196 pares (15,17\%). As proporções ficaram próximas da divisão planejada de 70\%, 15\% e 15\%, com pequenas diferenças esperadas quando a alocação é realizada por grupos, e não por registros individuais. A Tabela 10 mostra o particionamento da amostra.

\textbf{Tabela 10 -- Divisão agrupada da amostra}
\begin{longtable}[]{@{}
			>{\raggedright\arraybackslash}p{(\linewidth - 12\tabcolsep) * \real{0.1600}}
			>{\raggedright\arraybackslash}p{(\linewidth - 12\tabcolsep) * \real{0.0909}}
			>{\raggedright\arraybackslash}p{(\linewidth - 12\tabcolsep) * \real{0.1673}}
			>{\raggedright\arraybackslash}p{(\linewidth - 12\tabcolsep) * \real{0.1409}}
			>{\raggedright\arraybackslash}p{(\linewidth - 12\tabcolsep) * \real{0.1444}}
			>{\raggedright\arraybackslash}p{(\linewidth - 12\tabcolsep) * \real{0.1394}}
			>{\raggedright\arraybackslash}p{(\linewidth - 12\tabcolsep) * \real{0.1571}}@{}}
		\toprule\noalign{}
		\begin{minipage}[b]{\linewidth}\raggedright
			\textbf{Partição}
		\end{minipage} & \begin{minipage}[b]{\linewidth}\raggedright
			\textbf{Pares}
		\end{minipage} & \begin{minipage}[b]{\linewidth}\raggedright
			\textbf{\% dos pares}
		\end{minipage} & \begin{minipage}[b]{\linewidth}\raggedright
			\textit{\textbf{Commits}}
		\end{minipage} & \begin{minipage}[b]{\linewidth}\raggedright
			\textbf{Negativos}
		\end{minipage} & \begin{minipage}[b]{\linewidth}\raggedright
			\textbf{Positivos}
		\end{minipage} & \begin{minipage}[b]{\linewidth}\raggedright
			\textbf{Prevalência positiva}
		\end{minipage} \\
		\midrule\noalign{}
		\endhead
		\bottomrule\noalign{}
		\endlastfoot
		Treinamento & 902 & 69,81\% & 898 & 128 & 774 & 85,81\% \\
		Validação & 194 & 15,02\% & 192 & 24 & 170 & 87,63\% \\
		Teste & 196 & 15,17\% & 193 & 28 & 168 & 85,71\% \\
		Total & 1.292 & 100,00\% & --- & 180 & 1.112 & 86,07\% \\
\end{longtable}

A prevalência da classe positiva foi semelhante nas três partições: 85,81\% no treinamento, 87,63\% na validação e 85,71\% no teste. A amplitude entre as prevalências extremas foi de apenas 1,92 ponto percentual, indicando que a divisão agrupada preservou adequadamente a composição da variável de resposta. Ao mesmo tempo, a prevalência elevada evidencia desbalanceamento de classes: aproximadamente seis em cada sete pares foram classificados como integrados. Por essa razão, a avaliação posterior não deve se apoiar apenas na acurácia; métricas como precisão, revocação, área sob a curva precisão-revocação, \textit{Brier Score} e \textit{log loss} são mais informativas para verificar discriminação e qualidade probabilística.

\subsection{\texorpdfstring{Multicolinearidade no conjunto de
			treinamento}{Multicolinearidade no conjunto de treinamento}}\label{multicolinearidade-no-conjunto-de-treinamento}

Após a retirada dos preditores CoincideAssunto, Intervalo e CossenoDiffs, o diagnóstico de multicolinearidade foi repetido exclusivamente no conjunto de treinamento. A maior correlação remanescente ocorreu entre CoincideAutor e SimilaridadeDiffs (r=0,8744). As demais correlações foram inferiores a 0,40 em módulo, com destaque para CoincideAutor e CoincideIndices (r=0,3704), CoincideIndices e SimilaridadeDiffs (r = 0,3397) e CoincideLink e CoincideAutor (r=0,3185). A Tabela 11 apresenta a matriz de correlação das variáveis independentes do modelo reduzido.

\textbf{Tabela 11 -- Matriz de correlação dos preditores no conjunto de treinamento}
\begin{longtable}[]{@{}
			>{\raggedright\arraybackslash}p{(\linewidth - 10\tabcolsep) * \real{0.1830}}
			>{\raggedright\arraybackslash}p{(\linewidth - 10\tabcolsep) * \real{0.1622}}
			>{\raggedright\arraybackslash}p{(\linewidth - 10\tabcolsep) * \real{0.1695}}
			>{\raggedright\arraybackslash}p{(\linewidth - 10\tabcolsep) * \real{0.1768}}
			>{\raggedright\arraybackslash}p{(\linewidth - 10\tabcolsep) * \real{0.1169}}
			>{\raggedright\arraybackslash}p{(\linewidth - 10\tabcolsep) * \real{0.1917}}@{}}
		\toprule\noalign{}
		\begin{minipage}[b]{\linewidth}\raggedright
			\textbf{Variável}
		\end{minipage} & \begin{minipage}[b]{\linewidth}\raggedright
			\textbf{CoincideLink}
		\end{minipage} & \begin{minipage}[b]{\linewidth}\raggedright
			\textbf{CoincideAutor}
		\end{minipage} & \begin{minipage}[b]{\linewidth}\raggedright
			\textbf{CoincideIndices}
		\end{minipage} & \begin{minipage}[b]{\linewidth}\raggedright
			\textbf{Delta}
		\end{minipage} & \begin{minipage}[b]{\linewidth}\raggedright
			\textbf{SimilaridadeDiffs}
		\end{minipage} \\
		\midrule\noalign{}
		\endhead
		\bottomrule\noalign{}
		\endlastfoot
		CoincideLink & \textit{1,0000} & 0,3185 & 0,2181 & 0,0546 & 0,3063 \\
		CoincideAutor & 0,3185 & \textit{1,0000} & 0,3704 & 0,2055 & 0,8744 \\
		CoincideIndices & 0,2181 & 0,3704 & \textit{1,0000} & -0,0772 & 0,3397 \\
		Delta & 0,0546 & 0,2055 & -0,0772 & \textit{1,0000} & 0,1819 \\
		SimilaridadeDiffs & 0,3063 & 0,8744 & 0,3397 & 0,1819 & \textit{1,0000} \\
\end{longtable}

Embora a correlação de 0,8744 ainda indique forte associação bivariada entre CoincideAutor e SimilaridadeDiffs, os respectivos VIFs diminuíram para 4,4662 e 4,2686 como mostra os dados na Tabela 12. A redução é expressiva em comparação com o modelo inicial, no qual CoincideAutor apresentava VIF de 22,7139 e SimilaridadeDiffs VIF de 8,1359. Os demais preditores apresentaram VIFs próximos de 1. Esse padrão indica que a retirada das variáveis redundantes reduziu substancialmente a inflação da variância e permitiu manter, no modelo, dimensões conceitualmente distintas de autoria e similaridade textual.

\textbf{Tabela 12 -- Fator de inflação da variância no conjunto de
	treinamento}
\begin{longtable}[]{@{}
			>{\raggedright\arraybackslash}p{(\linewidth - 6\tabcolsep) * \real{0.2196}}
			>{\centering\arraybackslash}p{(\linewidth - 6\tabcolsep) * \real{0.1258}}
			>{\centering\arraybackslash}p{(\linewidth - 6\tabcolsep) * \real{0.1572}}
			>{\centering\arraybackslash}p{(\linewidth - 6\tabcolsep) * \real{0.4974}}@{}}
		\toprule\noalign{}
		\begin{minipage}[b]{\linewidth}\raggedright
			\textbf{Variável}
		\end{minipage} & \begin{minipage}[b]{\linewidth}\centering
			\textbf{VIF}
		\end{minipage} & \begin{minipage}[b]{\linewidth}\centering
			\textbf{Tolerância}
		\end{minipage} & \begin{minipage}[b]{\linewidth}\centering
			\textbf{Interpretação}
		\end{minipage} \\
		\midrule\noalign{}
		\endhead
		\bottomrule\noalign{}
		\endlastfoot
		CoincideAutor & 4,4662 & 0,2239 & Atenção, mas abaixo do limite
		operacional de 5 \\
		SimilaridadeDiffs & 4,2686 & 0,2343 & Atenção, mas abaixo do limite
		operacional de 5 \\
		CoincideIndices & 1,2095 & 0,8268 & Baixa multicolinearidade \\
		CoincideLink & 1,1312 & 0,8840 & Baixa multicolinearidade \\
		Delta & 1,0748 & 0,9304 & Baixa multicolinearidade \\
\end{longtable}

Os valores de VIF devem ser interpretados de forma contextual. O'Brien (2007) adverte que regras fixas não justificam, isoladamente, a exclusão de uma variável. Neste estudo, a permanência de CoincideAutor e
SimilaridadeDiffs é defensável porque ambas representam mecanismos substantivamente distintos, os VIFs ficaram abaixo de 5 no modelo reduzido e os coeficientes passaram a apresentar contribuição estatística no ajuste de treinamento. Ainda assim, a correlação remanescente recomenda cautela ao interpretar cada coeficiente como efeito isolado e reforça a necessidade de avaliar o desempenho em validação e teste.

\subsection{\texorpdfstring{Resultados do modelo logístico no
			conjunto de
			treinamento}{Resultados do modelo logístico no conjunto de treinamento}}\label{resultados-do-modelo-loguxedstico-no-conjunto-de-treinamento}

Uma vez concluído o diagnóstico, o modelo logístico-base foi ajustado às 902 observações de treinamento com os cinco preditores selecionados. O ajuste convergiu e o teste da razão de verossimilhanças foi globalmente
significativo (p \textless{} 0,001), indicando desempenho superior ao modelo com apenas o intercepto. O pseudo-R² de McFadden foi 0,8265 e a log-verossimilhança, -63,910. Embora o pseudo-R² elevado indique forte capacidade explicativa no treinamento, ele não representa proporção de variância explicada no mesmo sentido do R² da regressão linear. A generalização, portanto, depende dos resultados de validação e teste. As Tabelas 13 e 14 resumem o ajuste e os coeficientes.

\textbf{Tabela 13 -- Ajuste do modelo logístico-base no treinamento}
\begin{longtable}[]{@{}
			>{\raggedright\arraybackslash}p{(\linewidth - 2\tabcolsep) * \real{0.5035}}
			>{\raggedright\arraybackslash}p{(\linewidth - 2\tabcolsep) * \real{0.2821}}@{}}
		\toprule\noalign{}
		\begin{minipage}[b]{\linewidth}\raggedright
			\textbf{Indicador}
		\end{minipage} & \begin{minipage}[b]{\linewidth}\raggedright
			\textbf{Resultado}
		\end{minipage} \\
		\midrule\noalign{}
		\endhead
		\bottomrule\noalign{}
		\endlastfoot
		Observações & 902 \\
		Preditores & 5 \\
		Graus de liberdade residuais & 896 \\
		Log-verossimilhança & -63,910 \\
		Log-verossimilhança do modelo nulo & -368,390 \\
		Pseudo-R² de McFadden & 0,8265 \\
		p do teste da razão de verossimilhanças & \textless{} 0,001 \\
		Convergência & Sim \\
\end{longtable}

\newpage

\textbf{Tabela 14 -- Coeficientes do modelo logístico no treinamento}
\begin{longtable}[]{@{}
			>{\raggedright\arraybackslash}p{(\linewidth - 8\tabcolsep) * \real{0.2134}}
			>{\centering\arraybackslash}p{(\linewidth - 8\tabcolsep) * \real{0.1142}}
			>{\centering\arraybackslash}p{(\linewidth - 8\tabcolsep) * \real{0.1724}}
			>{\centering\arraybackslash}p{(\linewidth - 8\tabcolsep) * \real{0.1407}}
			>{\centering\arraybackslash}p{(\linewidth - 8\tabcolsep) * \real{0.3593}}@{}}
		\toprule\noalign{}
		\begin{minipage}[b]{\linewidth}\centering
			\textbf{Variável}
		\end{minipage} & \begin{minipage}[b]{\linewidth}\centering
			\textbf{$\beta$}
		\end{minipage} & \begin{minipage}[b]{\linewidth}\centering
			\textbf{Erro-padrão}
		\end{minipage} & \begin{minipage}[b]{\linewidth}\centering
			\textbf{p-value}
		\end{minipage} & \begin{minipage}[b]{\linewidth}\centering
			\textbf{IC 95\% de $\beta$}
		\end{minipage} \\
		\midrule\noalign{}
		\endhead
		\bottomrule\noalign{}
		\endlastfoot
		Intercepto & -5,3758 & 1,210 & \textless{} 0,001 & -7,747 a -3,005 \\
		CoincideLink & 3,1999 & 1,292 & 0,013 & 0,667 a 5,733 \\
		CoincideAutor & 6,4489 & 1,296 & \textless{} 0,001 & 3,909 a 8,988 \\
		CoincideIndices & 1,3705 & 0,665 & 0,039 & 0,067 a 2,674 \\
		Delta & -0,0650 & 0,171 & 0,703 & -0,399 a 0,269 \\
		SimilaridadeDiffs & 1,9876 & 0,664 & 0,003 & 0,687 a 3,288 \\
\end{longtable}


A análise dos coeficientes reforça a relevância das variáveis mantidas. CoincideAutor apresentou o maior coeficiente positivo entre os preditores binários ($\beta$=6,4489; p \textless{} 0,001), com intervalo de confiança integralmente acima de zero (3,909 a 8,988). Esse resultado fornece evidência forte de associação positiva entre coincidência de autoria e integração, mantidas constantes as demais variáveis.

CoincideLink também apresentou associação positiva e estatisticamente significativa ($\beta$=3,1999; p-value=0,013; IC 95\%: 0,667 a 5,733). Já CoincideIndices apresentou coeficiente positivo de 1,3705 e p = 0,039. Portanto, a coincidência entre os índices do diff acrescentou informação ao modelo após o controle pelos demais preditores, embora a evidência tenha sido menos forte do que a observada para CoincideAutor e SimilaridadeDiffs.

SimilaridadeDiffs permaneceu estatisticamente significativa depois da retirada de CossenoDiffs ($\beta$=1,9876; p-value=0,003). O intervalo de confiança, de 0,687 a 3,288, também permaneceu acima de zero. Assim, valores mais elevados de similaridade textual estiveram associados a maior probabilidade de integração, mantidos constantes os demais preditores.

Em contraste, Delta não apresentou associação estatisticamente identificável no treinamento ($\beta$=-0,0650; p-value=0,703; IC 95\%: -0,399 a 0,269). Como a variável foi padronizada, o coeficiente representa a alteração no preditor linear associada ao aumento de um desvio-padrão em Delta. Sua permanência no modelo final deve, portanto, apoiar-se em justificativa teórica ou em eventual ganho de desempenho nos conjuntos de validação e teste.

Por fim, o intercepto foi negativo e estatisticamente significativo ($\beta$=-5,3758; p-value \textless{} 0,001). Esse parâmetro representa o logaritmo das chances de integração quando todos os preditores assumem valor zero.

\section{Calibração probabilística e desempenho do modelo nos conjuntos de validação e teste}

Após o ajuste no treinamento, o desempenho foi examinado em quatro cenários complementares: probabilidades brutas e calibradas por  Platt (\cite{niculescu2005probabilities}) nos conjuntos de validação e teste. Essa comparação permite separar o ajuste do calibrador de sua capacidade de generalização.

\begin{itemize}
	\item
	\textbf{Validação - probabilidade bruta}: desempenho, no conjunto de
	validação, das probabilidades produzidas diretamente pelo modelo de
	regressão logística ajustado exclusivamente com os dados de
	treinamento (modelo-base), antes da calibração.
	\item
	\textbf{Validação - Platt}: desempenho, no conjunto de validação, das
	probabilidades produzidas pelo modelo ajustado no treinamento e
	posteriormente transformadas pelo calibrador de Platt. Como os
	parâmetros do calibrador foram estimados nesse mesmo conjunto, esses
	resultados descrevem seu ajuste, mas não constituem uma avaliação
	independente da capacidade de generalização da calibração.
	\item
	\textbf{Teste - probabilidade bruta}: desempenho, no conjunto de
	teste, das probabilidades produzidas diretamente pelo modelo de
	regressão logística ajustado exclusivamente com os dados de
	treinamento, sem aplicação do calibrador de Platt.
	\item
	\textbf{Teste - Platt}: desempenho, no conjunto de teste, das
	probabilidades inicialmente produzidas pelo modelo de regressão
	logística ajustado no treinamento e posteriormente transformadas pelo
	calibrador de Platt estimado no conjunto de validação. Essa avaliação
	mostra se a recalibração se generaliza para dados que não foram
	utilizados no ajuste do modelo-base nem do calibrador.
\end{itemize}

Inicialmente, são mostrados o tamanho e a prevalência dos conjuntos de dados usados em cada avaliação na Tabela 15.

\textbf{Tabela 15 -- Tamanho e prevalência positiva por avaliação}
\begin{longtable}[]{@{}
			>{\raggedright\arraybackslash}p{(\linewidth - 4\tabcolsep) * \real{0.4250}}
			>{\centering\arraybackslash}p{(\linewidth - 4\tabcolsep) * \real{0.1573}}
			>{\centering\arraybackslash}p{(\linewidth - 4\tabcolsep) * \real{0.2518}}@{}}
		\toprule\noalign{}
		\begin{minipage}[b]{\linewidth}\raggedright
			\textbf{Avaliação}
		\end{minipage} & \begin{minipage}[b]{\linewidth}\centering
			\textbf{n}
		\end{minipage} & \begin{minipage}[b]{\linewidth}\centering
			\textbf{Prevalência}
		\end{minipage} \\
		\midrule\noalign{}
		\endhead
		\bottomrule\noalign{}
		\endlastfoot
		Validação - probabilidade bruta & 194 & 87,63\% \\
		Validação - Platt & 194 & 87,63\% \\
		Teste - probabilidade bruta & 196 & 85,71\% \\
		Teste - Platt & 196 & 85,71\% \\
\end{longtable}

A prevalência de casos positivos foi semelhante entre validação (87,63\%) e teste (85,71\%), diferença de 1,92 ponto percentual. O predomínio da classe positiva exige cautela com métricas dependentes de limiar e torna especialmente relevante a área sob a curva precisão-revocação. Antes dessas métricas de classificação, porém, é
necessário examinar a qualidade das probabilidades, sintetizada na Tabela 16.

\textbf{Tabela 16 -- Qualidade das previsões probabilísticas}
\begin{longtable}[]{@{}
			>{\raggedright\arraybackslash}p{(\linewidth - 6\tabcolsep) * \real{0.3768}}
			>{\raggedright\arraybackslash}p{(\linewidth - 6\tabcolsep) * \real{0.2044}}
			>{\raggedright\arraybackslash}p{(\linewidth - 6\tabcolsep) * \real{0.2831}}
			>{\raggedright\arraybackslash}p{(\linewidth - 6\tabcolsep) * \real{0.1358}}@{}}
		\toprule\noalign{}
		\begin{minipage}[b]{\linewidth}\raggedright
			\textbf{Avaliação}
		\end{minipage} & \begin{minipage}[b]{\linewidth}\raggedright
			\textbf{Brier Score}
		\end{minipage} & \begin{minipage}[b]{\linewidth}\raggedright
			\textbf{Brier de referência}
		\end{minipage} & \begin{minipage}[b]{\linewidth}\raggedright
			\textbf{Log loss}
		\end{minipage} \\
		\midrule\noalign{}
		\endhead
		\bottomrule\noalign{}
		\endlastfoot
		Validação - probabilidade bruta & 0,0164 & 0,1087 & 0,0759 \\
		Validação - Platt & 0,0160 & 0,1087 & 0,0756 \\
		Teste - probabilidade bruta & 0,0222 & 0,1225 & 0,0892 \\
		Teste - Platt & 0,0212 & 0,1225 & 0,0843 \\
\end{longtable}

O \textit{Brier Score} corresponde ao erro quadrático médio entre a probabilidade prevista e o resultado observado; quanto menor, melhor. Na validação, o \textit{Platt scaling} reduziu o \textit{Brier Score} de 0,0164 para 0,0160, melhora absoluta de 0,000404, ou aproximadamente 2,47\%. No teste, a redução foi de 0,0222 para 0,0212, melhora absoluta de 0,001001, ou aproximadamente 4,51\%. A melhora no teste indica que a recalibração produziu probabilidades globalmente mais próximas dos resultados observados.

O \textit{Brier} de referência corresponde à previsão que atribui a todos os casos a prevalência positiva do treinamento. Seus valores, 0,1087 na validação e 0,1225 no teste, são muito superiores aos do modelo. No
teste, o \textit{Brier} calibrado foi cerca de 82,7\% menor que o da referência, evidenciando ganho probabilístico substancial em relação a uma previsão constante.

O \textit{log loss}, ou perda logarítmica, penaliza previsões erradas com maior intensidade quando são feitas com alta confiança. Quanto menor, melhor. Na validação, a redução após Platt foi pequena, de 0,0759 para 0,0756. No teste, caiu de 0,0892 para 0,0843, melhora aproximada de 5,52\%. Esse resultado favorece o \textit{Platt scaling}, pois sugere menor penalização por previsões excessivamente confiantes e incorretas
fora da amostra. Para complementar essa análise global, a Tabela 17 apresenta métricas específicas de calibração.

\textbf{Tabela 17 -- Métricas de calibração probabilística}
\begin{longtable}[]{@{}
			>{\raggedright\arraybackslash}p{(\linewidth - 6\tabcolsep) * \real{0.3925}}
			>{\raggedright\arraybackslash}p{(\linewidth - 6\tabcolsep) * \real{0.2516}}
			>{\raggedright\arraybackslash}p{(\linewidth - 6\tabcolsep) * \real{0.1729}}
			>{\raggedright\arraybackslash}p{(\linewidth - 6\tabcolsep) * \real{0.1830}}@{}}
		\toprule\noalign{}
		\begin{minipage}[b]{\linewidth}\raggedright
			\textbf{Avaliação}
		\end{minipage} & \begin{minipage}[b]{\linewidth}\raggedright
			\textbf{ECE (10 faixas)}
		\end{minipage} & \begin{minipage}[b]{\linewidth}\raggedright
			\textbf{Intercepto}
		\end{minipage} & \begin{minipage}[b]{\linewidth}\raggedright
			\textbf{Inclinação}
		\end{minipage} \\
		\midrule\noalign{}
		\endhead
		\bottomrule\noalign{}
		\endlastfoot
		Validação - probabilidade bruta & 0,007432 & 0,2930 & 0,9509 \\
		Validação - Platt & 0,008365 & -0,1201 & 1,0491 \\
		Teste - probabilidade bruta & 0,014887 & 1,0161 & 0,8180 \\
		Teste - Platt & 0,018041 & 0,6608 & 0,9024 \\
\end{longtable}

O ECE de dez faixas resume a diferença média ponderada entre aprobabilidade prevista e a frequência observada nos dez intervalos deprobabilidade. O ideal é zero. Ao contrário do \textit{Brier Score} e do \textit{log loss} o ECE aumentou após a calibração: de 0,007432 para 0,008365 na validação e de 0,014887 para 0,018041 no teste. No teste, o aumento foi de aproximadamente 21,2\%. Esse resultado indica que, para a divisão específica em dez faixas, a correspondência local entre previsões médias e frequências observadas piorou. O ECE deve ser interpretado com cautela porque depende dos limites das faixas e pode ser instável quando há poucas observações em algumas delas.

O intercepto de calibração ideal é zero. Valores positivos indicam que, depois de considerar a inclinação, o evento ocorreu mais frequentemente do que as probabilidades sugeriam; valores negativos apontam o padrão
inverso. No teste, o \textit{Platt scaling} aproximou o intercepto de zero, de 1,0161 para 0,6608, reduzindo a subestimação global, embora sem eliminá-la.

A inclinação de calibração ideal é um. Valores menores que um indicam probabilidades excessivamente extremas, enquanto valores maiores que um indicam probabilidades pouco extremas. No teste, a inclinação melhorou de 0,8180 para 0,9024, aproximando-se do ideal. Na validação, passou de 0,9509 para 1,0491; ambos os valores estão praticamente à mesma distância de um. Em conjunto, intercepto e inclinação favorecem o \textit{Platt scaling} no teste, apesar da piora do ECE.

No estudo foi usado o limiar de 0,50 na transformação das probabilidades de integração e não integração nas respectivas classes. Probabilidades maiores ou iguais a 0,50 foram classificadas como Sim (patch integrado), e probabilidades menores que 0,50 como Não (patch não integrado). Na Tabela 18 exibe as métricas de classificação do modelo logístico.

\textbf{Tabela 18 -- Precisão, revocação e F1 no limiar de 0,50}
\begin{longtable}[]{@{}
			>{\raggedright\arraybackslash}p{(\linewidth - 6\tabcolsep) * \real{0.4554}}
			>{\centering\arraybackslash}p{(\linewidth - 6\tabcolsep) * \real{0.1959}}
			>{\centering\arraybackslash}p{(\linewidth - 6\tabcolsep) * \real{0.1984}}
			>{\centering\arraybackslash}p{(\linewidth - 6\tabcolsep) * \real{0.1503}}@{}}
		\toprule\noalign{}
		\begin{minipage}[b]{\linewidth}\raggedright
			\textbf{Avaliação}
		\end{minipage} & \begin{minipage}[b]{\linewidth}\centering
			\textbf{Precisão\_050}
		\end{minipage} & \begin{minipage}[b]{\linewidth}\centering
			\textbf{Revocação\_050}
		\end{minipage} & \begin{minipage}[b]{\linewidth}\centering
			\textbf{F1\_050}
		\end{minipage} \\
		\midrule\noalign{}
		\endhead
		\bottomrule\noalign{}
		\endlastfoot
		Validação - probabilidade bruta & 98,27\% & 100,00\% & 99,13\% \\
		Validação - Platt & 98,27\% & 100,00\% & 99,13\% \\
		Teste - probabilidade bruta & 98,22\% & 98,81\% & 98,52\% \\
		Teste - Platt & 98,22\% & 98,81\% & 98,52\% \\
\end{longtable}

Precisão\_050 é a proporção de classificações positivas que realmente eram positivas. Revocação\_050 é a proporção dos positivos reais identificada pelo modelo. F1\_050 é a média harmônica entre precisão e revocação e resume o equilíbrio entre ambas. No teste, a precisão foi de 98,22\%, a revocação de 98,81\% e a F1 de 98,52\%. 

Os valores não mudaram após Platt, mostrando que nenhuma observação relevante atravessou o limiar de 0,50. Isso não significa que a calibração tenha sido inútil: ela pode melhorar o significado numérico das probabilidades sem alterar as classes nesse limiar.

Além do desempenho em um único limiar, avaliou-se a capacidade do modelo de ordenar os pares efetivamente relacionados acima dos não relacionados. A Tabela 19 mostra a área sob a curva precisão-revocação.

\textbf{Tabela 19 -- Área sob a curva precisão-revocação}
\begin{longtable}[]{@{}
			>{\raggedright\arraybackslash}p{(\linewidth - 2\tabcolsep) * \real{0.3768}}
			>{\raggedright\arraybackslash}p{(\linewidth - 2\tabcolsep) * \real{0.3144}}@{}}
		\toprule\noalign{}
		\begin{minipage}[b]{\linewidth}\raggedright
			\textbf{Avaliação}
		\end{minipage} & \begin{minipage}[b]{\linewidth}\raggedright
			\textbf{Área precisão-revocação}
		\end{minipage} \\
		\midrule\noalign{}
		\endhead
		\bottomrule\noalign{}
		\endlastfoot
		Validação - probabilidade bruta & 0,9941 \\
		Validação - Platt & 0,9941 \\
		Teste - probabilidade bruta & 0,9958 \\
		Teste - Platt & 0,9958 \\
\end{longtable}

A área sob a curva precisão-revocação resume a capacidade de ordenar os casos positivos acima dos negativos ao longo de todos os limiares. Valores próximos de um indicam discriminação excelente. O resultado foi 0,9941 na validação e 0,9958 no teste. A área permaneceu inalterada após Platt porque a calibração aplica uma transformação monotônica às probabilidades: seus valores mudam, mas a ordenação dos registros tende a ser preservada. A Tabela 20 exibe o quadro resumo do processo de calibração.

\textbf{Tabela 20 -- Efeito da calibração de Platt}
\begin{longtable}[]{@{}
			>{\raggedright\arraybackslash}p{(\linewidth - 6\tabcolsep) * \real{0.2482}}
			>{\raggedright\arraybackslash}p{(\linewidth - 6\tabcolsep) * \real{0.2490}}
			>{\raggedright\arraybackslash}p{(\linewidth - 6\tabcolsep) * \real{0.1940}}
			>{\raggedright\arraybackslash}p{(\linewidth - 6\tabcolsep) * \real{0.3088}}@{}}
		\toprule\noalign{}
		\begin{minipage}[b]{\linewidth}\raggedright
			\textbf{Critério}
		\end{minipage} & \begin{minipage}[b]{\linewidth}\raggedright
			\textbf{Validação}
		\end{minipage} & \begin{minipage}[b]{\linewidth}\raggedright
			\textbf{Teste}
		\end{minipage} & \begin{minipage}[b]{\linewidth}\raggedright
			\textbf{Interpretação}
		\end{minipage} \\
		\midrule\noalign{}
		\endhead
		\bottomrule\noalign{}
		\endlastfoot
		Brier \textit{Score} & Melhora pequena & Melhora & Favorável ao Platt \\
		\textit{Log loss} & Melhora mínima & Melhora & Favorável ao Platt \\
		ECE em 10 faixas & Piora & Piora & Desfavorável ao Platt \\
		Intercepto & Aproxima de zero & Aproxima de zero & Favorável ao Platt \\
		Inclinação & Permanece próxima de 1 & Aproxima de 1 & Favorável ao
		Platt \\
		Precisão, revocação e F1 em 0,50 & Sem alteração & Sem alteração &
		Neutro \\
		Área precisão-revocação & Sem alteração & Sem alteração & Esperado em
		transformação monotônica \\
\end{longtable}

Considerando em conjunto, a evidência é moderadamente favorável ao Platt \textit{scaling}. No conjunto de teste, houve melhora do \textit{Brier Score,} do \textit{log loss}, do intercepto e da inclinação de calibração, sem perda na classificação ao limiar de 0,50 nem na discriminação. A principal evidência contrária foi o aumento do ECE de dez faixas. Como \textit{Brier Score} e \textit{log loss} são métricas próprias de avaliação probabilística e consideram cada observação individualmente, a melhora de ambas no teste fornece um argumento mais robusto em favor das probabilidades calibradas.

Vale ressaltar que o modelo já apresentava desempenho muito elevado antes da calibração. Portanto, o ganho de Platt foi incremental, não transformador. Para decisões binárias com limiar 0,50, as versões bruta e calibrada são equivalentes nesta amostra. Para usos em que o valor da probabilidade é relevante, como categorização por níveis de confiança, priorização de revisão manual ou comparação de risco, a versão \textit{Platt} é preferível. A Tabela 21 encerra a seção com um guia de interpretação das métricas utilizadas.

\textbf{Tabela 21 -- Guia de interpretação}
\begin{longtable}[]{@{}
			>{\raggedright\arraybackslash}p{(\linewidth - 4\tabcolsep) * \real{0.2982}}
			>{\raggedright\arraybackslash}p{(\linewidth - 4\tabcolsep) * \real{0.3724}}
			>{\raggedright\arraybackslash}p{(\linewidth - 4\tabcolsep) * \real{0.3293}}@{}}
		\toprule\noalign{}
		\begin{minipage}[b]{\linewidth}\raggedright
			\textbf{Métrica}
		\end{minipage} & \begin{minipage}[b]{\linewidth}\raggedright
			\textbf{O que avalia}
		\end{minipage} & \begin{minipage}[b]{\linewidth}\raggedright
			\textbf{Valor desejável}
		\end{minipage} \\
		\midrule\noalign{}
		\endhead
		\bottomrule\noalign{}
		\endlastfoot
		Prevalência & Proporção observada de positivos & Descritivo; não há
		valor ideal \\
		Brier Score & Erro quadrático das probabilidades & Menor \\
		Brier de referência & Erro de prever a prevalência do treino para todos
		& Modelo deve ser menor \\
		ECE de 10 faixas & Diferença entre previsão média e frequência observada
		por faixa & Próximo de 0 \\
		Log loss & Penalização probabilística, forte para erros confiantes &
		Menor \\
		Intercepto de calibração & Desvio global das probabilidades & Próximo de
		0 \\
		Inclinação de calibração & Extremidade ou dispersão das probabilidades &
		Próximo de 1 \\
		Precisão\_050 & Correção entre os positivos previstos no corte 0,50 &
		Maior \\
		Revocação\_050 & Cobertura dos positivos reais no corte 0,50 & Maior \\
		F1\_050 & Equilíbrio entre precisão e revocação no corte 0,50 & Maior \\
		Área precisão-revocação & Ordenação de positivos e negativos em todos os
		cortes & Próximo de 1 \\
\end{longtable}

